\documentclass[11pt]{article}
\usepackage[margin=1in]{geometry}
\usepackage{amsmath,amssymb,amsthm}
\usepackage{booktabs}
\usepackage{listings}
\usepackage{xcolor}
\usepackage[numbers]{natbib}
\usepackage{hyperref}

\lstdefinelanguage{lean}{
  morekeywords={def,theorem,lemma,example,namespace,end,import,where,instance,fun,by,decide,Prop,Type,open,intro,exact,rw,refine,have,simpa,apply,with,induction,using},
  sensitive=true, morecomment=[l]{--}, morecomment=[s]{/-}{-/}, morestring=[b]"
}
\lstset{language=lean, basicstyle=\ttfamily\small, breaklines=true, columns=fullflexible,
  keywordstyle=\color{blue!60!black}\bfseries, commentstyle=\color{gray},
  literate={∀}{{$\forall$}}1 {→}{{$\to$}}1 {↔}{{$\leftrightarrow$}}1 {∧}{{$\land$}}1
    {∨}{{$\lor$}}1 {¬}{{$\neg$}}1 {∈}{{$\in$}}1 {≤}{{$\le$}}1 {∏}{{$\prod$}}1
    {ℝ}{{$\mathbb{R}$}}1 {ℕ}{{$\mathbb{N}$}}1 {ι}{{$\iota$}}1 {·}{{$\cdot$}}1 {−}{{$-$}}1
    {⊤}{{$\top$}}1 {⊥}{{$\bot$}}1 {⟨}{{$\langle$}}1 {⟩}{{$\rangle$}}1 {←}{{$\leftarrow$}}1
    {⦃}{{$\{$}}1 {⦄}{{$\}$}}1}

\newtheorem{theorem}{Theorem}
\newtheorem{lemma}[theorem]{Lemma}
\newtheorem{proposal}{Proposed statement}

\title{Trees are the minimal connected corridor networks,\\ and therefore have no spare corridor:\\ a note for corridor planners with a Lean proof}
\author{Justin Hart\\Viridis LLC}
\date{24 November 2026}

\begin{document}
\maketitle

\begin{abstract}
Hierarchical Dendritic Forest Management (HDFM) designs tree-shaped (dendritic) corridor networks
between habitat patches. We treat the sealed statement \texttt{P6.MinimalNetwork}, which models patches as the
vertices and corridors as the edges of a finite simple graph and has three parts.
(1) A connected network on $n$ patches has at least $n-1$ corridors.
(2) It is connected with exactly $n-1$ corridors if and only if it is a tree.
(3) In a tree, removing any one corridor disconnects the network.
Together these fix a trade-off for the planner. A dendritic network uses the fewest corridors that can
connect the patches, and for exactly that reason none of its corridors is spare. We give a Lean~4 proof
of the statement against the pinned Mathlib, built from existing Mathlib lemmas plus a short bridge between
Mathlib's \texttt{Nat.card} counts and the statement's \texttt{Fintype} counts. These are classical facts
of finite graph theory, not new theorems. The graph carries no geometry, width, quality or cost, so the
result says nothing about whether any mapped corridor is used by any species. It does not discharge the
FaceCount assumption of \emph{Trees Cannot Cut} (project P1).
\end{abstract}

\section{Introduction}
Graph models of landscape connectivity represent habitat patches as vertices and functional
links between them as edges~\citep{urbankeitt2001}. Urban and Keitt argued that the minimum spanning
tree is a useful guide to which patches and links matter most to connectivity~\citep{urbankeitt2001}.
HDFM goes further and \emph{designs} dendritic corridor networks. A planner then needs two facts in
precise form: how few corridors suffice to connect $n$ patches, and what is given up by using that few.

Project P6 of Viridis Research builds the formal core of a corridor-network kernel. The first session
(P6-T001, 23 November 2026) proposed three statements. Justin Hart sealed the first,
\texttt{P6.MinimalNetwork}, by decision D-024. This session (P6-T002) proves it.

\paragraph{Question.} For a finite network of patches and corridors: is $n-1$ a lower bound on the
number of corridors of a connected network, which networks attain it, and do those networks have
any corridor whose loss they survive?

\paragraph{Answer.} Yes, $n-1$ is a lower bound. It is attained exactly by trees, and a tree survives
the loss of none of its corridors. Section~\ref{sec:tradeoff} spells out what this means for planning.

\section{Background and related work}
All three parts are textbook graph theory~\citep{diestel2017}. A connected graph on $n$ vertices has
at least $n-1$ edges, with equality exactly for trees. An edge is a bridge, meaning its removal
separates its endpoints, exactly when it lies on no cycle. Hence a graph is acyclic if and only if
every edge is a bridge. We use Lean~4~\citep{lean4} and Mathlib~\citep{mathlib2020}. Mathlib already
contains the needed facts, stated with \texttt{Nat.card}:
\texttt{SimpleGraph.Connected.card\_vert\_le\_card\_edgeSet\_add\_one},
\texttt{SimpleGraph.isTree\_iff\_connected\_and\_card},
\texttt{SimpleGraph.isAcyclic\_iff\_forall\_edge\_isBridge} and \texttt{SimpleGraph.isBridge\_iff}.
The contribution of this note is not the mathematics. It is a sealed statement whose hypotheses are listed
in full, a machine-checkable proof of it, and a plain account of what that proof does and does not cover.

\section{The model and the statement}
Patches form a finite type \texttt{V : Type} with decidable equality. Corridors are the edges of a
\texttt{G : SimpleGraph V} with decidable adjacency. The module writes $n=\texttt{patchCount V}=\texttt{Fintype.card V}$ and
$m=\texttt{corridorCount G}=\texttt{G.edgeFinset.card}$. Removing corridor $e$ is \texttt{G.deleteEdges \{e\}}. A simple
graph has no loops and at most one corridor between two patches. The model has no geometry, width,
quality or cost; a corridor either exists or does not.

\begin{theorem}[\texttt{P6.MinimalNetwork}]\label{thm:main}
For every finite $V$ and every simple graph $G$ on $V$:
\begin{enumerate}
\item if $G$ is connected then $n\le m+1$;
\item $G$ is connected and $m+1=n$ if and only if $G$ is a tree;
\item if $G$ is a tree then, for every edge $e$ of $G$, the graph $G-e$ is not connected.
\end{enumerate}
\end{theorem}

The sealed Lean definition is:
\begin{lstlisting}
def Statement : Prop :=
  ∀ (V : Type) [Fintype V] [DecidableEq V] (G : SimpleGraph V) [DecidableRel G.Adj],
    (G.Connected → patchCount V ≤ corridorCount G + 1) ∧
    (G.Connected ∧ corridorCount G + 1 = patchCount V ↔ G.IsTree) ∧
    (G.IsTree → ∀ e ∈ G.edgeSet, ¬ (G.deleteEdges {e}).Connected)
\end{lstlisting}
The instances are quantified inside \texttt{Statement}, so the proof must work for any
\texttt{Fintype} and decidability instances, not just the canonical ones.

\section{Proof}
\paragraph{Mathematical proof.} (1) Choose a spanning tree of the connected graph $G$. A tree on $n$
vertices has $n-1$ edges (induction, removing a leaf), so $G$ has at least $n-1$ edges.
(2) A tree is connected with $n-1$ edges by the same count. Conversely, suppose $G$ is connected with
$m=n-1$ and has a cycle. Deleting an edge $e$ of the cycle keeps $G$ connected, since the rest of the
cycle reconnects the endpoints of $e$. But $G-e$ has $n-2$ edges, which contradicts~(1). So $G$ is
acyclic, hence a tree.
(3) Let $G$ be a tree and $e=uv$ an edge. If $G-e$ were connected, a path from $u$ to $v$ in $G-e$
together with $e$ would form a cycle in $G$. So $u$ and $v$ are not joined in $G-e$, and $G-e$ is
disconnected.

\paragraph{Lean proof.} Parts (1) and (2) are Mathlib's lemmas once the counts are translated:
\texttt{Fintype.card V = Nat.card V} and \texttt{G.edgeFinset.card = Nat.card G.edgeSet}. Because
Mathlib's lemmas are stated with \texttt{Nat.card}, which does not depend on instances, this translation
also handles arbitrary \texttt{Fintype} instances. Part (3) uses acyclicity to make $e$ a bridge, splits
$e$ as \texttt{s(u, v)}, and notes that \texttt{deleteEdges} unfolds to the graph difference
\texttt{G \textbackslash{} fromEdgeSet \{s(u, v)\}} that appears in \texttt{isBridge\_iff}. The complete
candidate file follows.

\begin{lstlisting}
import Mathlib
import ViridisStatements.P6.MinimalNetwork

namespace Candidate

open SimpleGraph ViridisStatements.P6.MinimalNetwork

lemma corridorCount_eq {V : Type} [Fintype V] [DecidableEq V] (G : SimpleGraph V)
    [DecidableRel G.Adj] : corridorCount G = Nat.card G.edgeSet := by
  rw [corridorCount, Nat.card_eq_fintype_card, ← edgeFinset_card]

lemma patchCount_eq (V : Type) [Fintype V] : patchCount V = Nat.card V := by
  rw [patchCount, Nat.card_eq_fintype_card]

lemma tree_delete_not_connected {V : Type} (G : SimpleGraph V) (hT : G.IsTree)
    (e : Sym2 V) (he : e ∈ G.edgeSet) : ¬ (G.deleteEdges {e}).Connected := by
  have hb : G.IsBridge e := isAcyclic_iff_forall_edge_isBridge.mp hT.IsAcyclic he
  induction e using Sym2.ind with
  | _ u v =>
    rw [isBridge_iff] at hb
    intro hc
    apply hb.2
    have := hc.preconnected u v
    simpa [deleteEdges] using this

theorem main : ViridisStatements.P6.MinimalNetwork.Statement := by
  intro V _ _ G _
  refine ⟨?_, ?_, ?_⟩
  · intro hc
    rw [patchCount_eq, corridorCount_eq]
    exact hc.card_vert_le_card_edgeSet_add_one
  · rw [patchCount_eq, corridorCount_eq, isTree_iff_connected_and_card]
  · intro hT e he
    exact tree_delete_not_connected G hT e he

end Candidate
\end{lstlisting}

\section{Results}
The candidate \texttt{proofs/P6.MinimalNetwork.lean} compiles under Lean \texttt{v4.28.0} with Mathlib
commit \texttt{8f9d9cff6bd728b17a24e163c9402775d9e6a365}. It passed the local check script, which
runs the denylist, compile, statement-match, axiom and kernel-replay stages. Its only axioms are
\texttt{propext}, \texttt{Classical.choice} and \texttt{Quot.sound}. Whether the claim is labelled
proven is decided by the independent verifier, whose certification page precedes this paper. The sealed
module's own sanity checks still apply: the path on three patches attains $m+1=n$ and loses
connectivity when an edge is removed, and the triangle with an isolated fourth patch shows that the count
$m+1=n$ without connectivity does not make a tree.

\section{The trade-off for a corridor planner}\label{sec:tradeoff}
Read together, the three parts make one point. Part~1 is a budget floor: no plan connects $n$ patches
with fewer than $n-1$ corridors. Part~2 says the floor is reached exactly by dendritic (tree) plans, so
a planner who wants the fewest corridors must build a tree, and any connected plan with $n-1$
corridors is already one. Part~3 is the price: every corridor of such a plan is a bridge. Losing any single
corridor, through a road, a fire or a clearing, splits the network into two parts that can no longer reach each other
through the network. Redundancy costs corridors. On $n\ge 2$ patches, surviving any one loss needs at least $n$ corridors
(the proposed \texttt{P6.SingleFailureRobust}, task P6-T003, not proved here). The planner's choice between
the cheapest network and a robust one is therefore a choice about corridor count, not only about layout.

\section{Limitations}
\begin{itemize}
\item These are classical facts about finite graphs, not new theorems.
\item The graph carries no geometry, width, quality or cost. The result is about the \emph{number} of
corridors, not their length or price. Minimum-cost networks, vertex (patch) failures and simultaneous
multiple failures are not covered.
\item ``Connected'' means connected in the abstract graph. A proof certifies the mathematics of the model
only. It does not certify that any mapped corridor is used by any species, that a corridor that exists on a map
carries movement, or any ecological outcome.
\item Multigraphs (parallel corridors between the same two patches) are outside the model. With a doubled
corridor, part~3 would no longer describe what removing one physical corridor does.
\item It does not discharge the FaceCount assumption of \emph{Trees Cannot Cut} (project P1).
\end{itemize}

\section{Reproducibility}
The candidate proof is \texttt{proofs/P6.MinimalNetwork.lean}. It imports only \texttt{Mathlib} and the sealed
module \texttt{ViridisStatements.P6.MinimalNetwork} (sha256 \texttt{972f34a0\ldots 564b}) and uses none
of the forbidden constructs. To check it, run
\texttt{viridis-lean-check proofs/P6.MinimalNetwork.lean P6.MinimalNetwork} under toolchain
\texttt{leanprover/lean4:v4.28.0} and Mathlib commit \texttt{8f9d9cff6bd728b17a24e163c9402775d9e6a365}.
No data or numerical code are involved.

\section*{Changelog}
\begin{itemize}
\item 2026-11-23 (P6-T001): first version. Drafted three statement proposals
(\texttt{P6.MinimalNetwork}, \texttt{P6.SingleFailureRobust}, \texttt{P6.SerialParallel}) with sanity
checks, negative controls and rationales.
\item 2026-11-24 (P6-T002): \texttt{P6.MinimalNetwork} sealed (D-024). Rewrote the paper as a note
on that statement. Added the mathematical proof, a Lean proof, and a section on the planning trade-off.
Removed the unused relay-reliability reference. The two remaining proposals are summarised in the text and await their own tasks.
\end{itemize}

\bibliographystyle{plainnat}
\bibliography{refs}
\end{document}
